\section{Motivation}\label{sec:intro_motivation}

\subsection{Industry Motivation: BMW Quality Management}\label{sec:intro_motivation_bmw}

In modern automotive manufacturing, maintaining product quality requires the
rapid identification and resolution of defects throughout the production
process. When a quality incident is detected, one of the first and most
critical challenges faced by quality engineers is determining the
responsible causer of the defect: specifically, whether the issue
originates from a supplier-provided component or from an internal process
within production, assembly, logistics, or engineering. This distinction
directly impacts the speed and effectiveness of corrective actions. The
current BMW quality process follows this principle by forwarding
supplier-related incidents to the quality management trend desk, where a
quality notification is created in the SAP system and a corresponding
supplier complaint and 8D problem-solving process is initiated; internal
issues follow an alternative investigation and resolution process.

Despite extensive quality management procedures and the availability of
large amounts of historical quality data, accurately identifying the true
causer of a defect remains a challenging task. Quality associates often
rely on historical incidents, expert knowledge, rework records, inspection
information, supplier complaints, and previous 8D reports to support their
decisions. However, these data sources are distributed across different
systems and are rarely connected in a manner that supports efficient
cross-referencing. In practice, such approaches perform adequately when an
identical incident has occurred previously, but become significantly less
effective when new or slightly different defect scenarios arise. As a
result, valuable information contained in related incidents, rework
histories, inspections, and historical supplier interactions often remains
unused.

The consequences of incorrect causal attribution can be substantial. If a
defect is incorrectly classified as supplier-related, an unnecessary
supplier complaint and lengthy clarification cycle may be initiated.
Conversely, if a genuine supplier issue is not identified at an early
stage, corrective actions may be delayed while the underlying defects
continue to affect production. Furthermore, the formal confirmation of
supplier responsibility often depends on the outcome of subsequent
root-cause investigations, which may require considerable time before a
definitive acknowledgement becomes available. During this period, the
production organization must continue operating under uncertainty,
increasing the risk of repeated defects, additional rework activities, and
inefficient resource utilization. Reliable causer identification during
the earlier stages of defect analysis therefore represents a significant
opportunity for improving operational efficiency and reducing problem
resolution lead times.

At the same time, the organization possesses an enormous amount of
accumulated quality knowledge: historical quality incidents, rework
records, inspection records, supplier complaints, root causes, corrective
actions, and 8D reports collectively contain valuable information that
can support future decisions. However, this knowledge is often fragmented
across multiple enterprise systems and remains difficult to reuse
systematically. Consequently, quality engineers repeatedly perform similar
investigations even though evidence from comparable historical cases may
already exist. The challenge is therefore no longer the availability of
data, but rather the ability to transform historical enterprise knowledge
into actionable and trustworthy decision support.

To address this challenge, this thesis investigates the development of a
system capable of identifying likely defect causers based on evidence
derived from historical quality knowledge. The proposed solution
integrates quality incidents, rework activities, inspection information,
supplier-related reports, material data, and historical resolutions into
a unified knowledge graph. By utilizing this interconnected knowledge
base, the system aims to support quality engineers in determining whether
a defect is likely supplier-related or internally caused, identifying
recurring patterns, discovering lessons learned from previous incidents,
and providing transparent recommendations supported by historical
evidence. The objective is not to replace human expertise, but to augment
it through explainable and evidence-driven artificial intelligence.

\subsection{Scientific Motivation}\label{sec:intro_motivation_sci}

While the industry motivation of this work originates from automotive
quality management, the underlying research problem extends far beyond
the BMW domain. Similar challenges occur in many complex environments,
including aircraft maintenance, semiconductor manufacturing, healthcare
incident analysis, industrial fault diagnosis, and large-scale asset
management. In all of these settings, organizations possess vast amounts
of historical knowledge distributed across heterogeneous data systems,
and their decision-makers must continuously evaluate previously unseen
situations under uncertainty.

Traditional approaches for supporting such decisions generally fall into
three categories. Rule-based systems provide transparent decisions by
relying on predefined expert knowledge and deterministic rules; however,
their effectiveness is limited by their inability to generalize beyond
explicitly defined scenarios, and when no matching rule exists, the
system is unable to provide a meaningful recommendation. Machine learning
approaches improve predictive performance by learning patterns directly
from data, but they often operate as black boxes, making it difficult for
domain experts to understand why a specific recommendation was produced;
this lack of transparency limits trust and adoption in industrial
decision-making environments. Knowledge graph approaches provide a
structured representation of enterprise knowledge and support transparent
retrieval of related information through explicit relationships between
entities; however, most existing applications focus primarily on
deterministic retrieval rather than on reasoning over multiple
heterogeneous sources of evidence.

This thesis argues that none of these approaches alone is sufficient for
supporting complex industrial decision-making tasks. Instead, effective
decision support requires a combination of structured enterprise
knowledge, evidence retrieval, evidence evaluation, reasoning, and
explainability. Consequently, this research proposes an evidence-centric
reasoning framework for industrial knowledge graphs. Rather than producing
direct predictions from isolated data points, the framework retrieves
relevant historical evidence from multiple sources, evaluates the strength
and consistency of this evidence, resolves potential conflicts between
evidence channels, and generates recommendations through large language
model reasoning. The resulting recommendation is accompanied by a
traceable explanation that explicitly identifies the evidence supporting
the conclusion.

The core research challenge addressed by this thesis is therefore not
merely the prediction of a defect's causer. Instead, it investigates a
broader question:

\begin{quote}
\textbf{How can a system support trustworthy decision-making by
retrieving, aggregating, evaluating, and reasoning about heterogeneous
evidence contained within industrial knowledge graphs?}
\end{quote}

To answer this question, the proposed framework integrates symbolic models, graph
retrieval, multi-source evidence aggregation and ranking, conflict
resolution mechanisms, and large-language-model-based explanation
generation. Through this combination, the system seeks to overcome the
limitations of traditional rule-based systems, black-box machine learning
models, and purely deterministic knowledge-graph retrieval.

The scientific contribution of this work therefore extends beyond the
automotive domain. While BMW quality management provides the validation
environment, the proposed methodology represents a general framework for
explainable decision support on industrial knowledge graphs. The central
contribution lies in demonstrating how heterogeneous enterprise knowledge
graphs can be transformed into transparent, evidence-based recommendations
that remain understandable, auditable, and trustworthy for human
decision-makers.
